\section{Applying the convergence estimates to the selected histories}
\label{sec:selected-history-limits}
The preceding convergence arguments concern particular arrays, paths and
stationary branches. We now identify these objects along one construction.
This identification is needed because equality of the actions at one
period would not identify their coefficient arrays, and because a
cutoff-dependent change of the top boundary need not preserve a limit
unless it eventually leaves each fixed scale unchanged.

\subsection{One simultaneous stationary construction}
Fix the scale schedule, clipping data and coordinate conventions. Denote
the terminal cutoff by $N$ and the scale by $m\le N$. The finite
stationary construction of Section~\ref{sec:path-attachment-producer}
gives paths $R_{N,m}$, scalar corrections $\alpha_{N,m}$ and stationary
maps $\rho_{N,m}^{p,F}$, where $p$ is a period and $F$ an allowed forest.
All these objects are chosen together. Their one-step identities use
$R_{N,m+1}$ as input and $R_{N,m}$ as output. The path bounds hold at
the prescribed scale widths and allowances.

For the convergence assertion we use the strengthened version of this
construction. At scale $m$ it fixes three coefficient Hilbert spaces,
with encoding widths $\delta_{m,2},\delta_{m,3},\delta_{m,4}$ strictly
larger than the paid width $w_m$. The nonlinear path terms are encoded
in their canonical selected subspaces. It also supplies the following
uniformity: every admissible replacement of the input path with the
same attached row and with all path and loop allowances enlarged by
$64$ has an output in these same selected spaces, at the ordinary
output allowances. The output has the same classical and thermal
matching identities on every period and forest. Finally, the underlying
free horizon converges at each scale in all of its retained components,
and the root rows have an eventual bound at some width $v_m>w_m$.
These are the extended-domain, canonical-readout and retained-horizon
hypotheses used in Section~\ref{sec:background-path-convergence}.
They belong to the one-step analytic construction; neither finite
backward induction nor the argument below supplies them.

Here is why its coefficient representation is the one needed in that
section. Select the three coefficient arrays of each path once, at
$(N,m)$, and use that same selection in every occurrence of the path.
For $m+2<N$, the input arrays of the step at $m$ are then exactly the
output arrays of the step at $m+1$. The row at either occurrence is
also the same attached root row. Thus the input/output equality holds
in the coefficient spaces, before taking any limit. The one-step
identities and bounds are available for all these chosen arrays.
For a replacement of one array the other arrays remain fixed; their
closed linear constraints and their separate budgets are preserved.
The enlarged input domain therefore supplies exactly the additional
steps used in the contraction argument.

Apply Section~\ref{sec:background-path-convergence} to these data.
At every fixed scale there are one limiting path $R_m$ and one
scalar $\alpha_m$ such that every coefficient of $R_{N,m}$ converges
to the corresponding coefficient of $R_m$, and
$\alpha_{N,m}\to\alpha_m$. The limiting path has the original paid
bounds. For every fixed period, forest and derivative order $r$,
\begin{equation}\label{eq:selected-stationary-limit}
 D^r\rho_{N,m}^{p,F}(0)\longrightarrow D^r\rho_m^{p,F}(0).
\end{equation}
The limiting stationary germ is determined by $R_{m+1}$, its retained
quadratic data and its forest. Its Hessian is positive with the same
lower bound. Consequently no choice of a new subsequence for a forest
or derivative order is involved in \eqref{eq:selected-stationary-limit}.
The scalar limit is the one obtained from these same paths and thermal
identities.

\subsection{A single width for all four coefficient families}
The common dictionaries are part of the conclusion just obtained.
Write $r_{N,m}(a)$ and $r_m(a)$ for row coefficients on the canonical
row dictionary, and $c_{N,m,q}(a,\nu)$ and $c_{m,q}(a,\nu)$ for the
actual, exponentially encoded coefficients of the atoms of degree
$q=2,3,4$. Their supports, lengths $\ell(a)$ and color multi-indices
$\nu$ are independent of $N$ on a tail. Choose
\[
 w_m<\widetilde w_m<
 \min\{v_m,\delta_{m,2},\delta_{m,3},\delta_{m,4}\}.
\]
Then, for every $\varepsilon>0$, all sufficiently large $N$ satisfy
\begin{align}
 \sum_a |r_{N,m}(a)-r_m(a)|e^{\widetilde w_m\ell(a)}
       &\le\varepsilon,\label{eq:selected-row-width}\\
 \sum_{a,\nu}|c_{N,m,q}(a,\nu)-c_{m,q}(a,\nu)|
                         e^{\widetilde w_m\ell(a)}
       &\le\varepsilon\qquad(q=2,3,4).
 \label{eq:selected-atom-width}
\end{align}
For the row statement one may either sum the finitely many plaquette
directions or impose the estimate separately on each; shrinking the
requested $\varepsilon$ by their number gives the displayed version.

For completeness, the two tail bounds that permit this simultaneous
choice are different. If the row budget at $v_m$ is $B_m$, its
coefficient difference at the smaller width is bounded by
\[2B_m e^{-(v_m-\widetilde w_m)\ell(a)}.\]
For an atom, remove its
encoding factor $e^{-\delta_{m,q}\ell(a)}$. The resulting coefficient
arrays have a common Hilbert bound $C_{m,q}$, so their individual
differences have absolute value at most $2C_{m,q}$. Their weighted
majorant is therefore
$2C_{m,q}e^{-(\delta_{m,q}-\widetilde w_m)\ell(a)}$.
Both majorants are summable on the canonical dictionaries by the
geometric sums in Section~\ref{sec:fixed-rung-convergence}. The
coordinate limits and a finite-head, summable-tail argument prove
\eqref{eq:selected-row-width}--\eqref{eq:selected-atom-width}.
There are only four families and finitely many row directions, so
their eventual estimates hold on one tail. This gives one path
limit with one strictly larger width, as required in the local
integral comparison.

\subsection{Changing the terminal normalization}
The final Wilson normalization replaces the top path at $m=N$.
Let $\widetilde R_{N,m}$, $\widetilde\alpha_{N,m}$ and
$\widetilde\rho_{N,m}^{p,F}$ denote the resulting data. The change
of coordinates retains the exact identities
\begin{align*}
 \widetilde R_{N,m}&=R_{N,m}&& (m<N),\\
 \widetilde\alpha_{N,m}&=\alpha_{N,m}&&\text{for all }N,m,\\
 \widetilde\rho_{N,m}^{p,F}&=\rho_{N,m}^{p,F}&&(m+1<N).
\end{align*}
It also leaves the free bed unchanged for $m<N$ and the entire
associated tower unchanged for $m+1<N$; its transport row is
unchanged already for $m<N$. These are identities of the maps and
coefficient data, not merely equalities of their evaluated actions.

For fixed $m$, all the identities hold once $N>m+1$. Deleting this
finite initial segment changes none of the limits above, none of
the eventual larger-width bounds and none of the retained free
horizon limits. Equality of the stationary maps gives equality of
every derivative at zero. Thus the final normalized construction
has exactly the same $R_m$, $\alpha_m$ and stationary jets. In
particular, replacing the terminal path has not introduced a second
background for the activity estimates.

\subsection{The activity and coupling limits on this background}
Fix a period $p$. The regular supports, forced inventories and norms
at scale $m$ depend on the fixed scale and period, not on $N$.
Embed their actual arrays into the one completion $\mathcal H_{m,p}$
of Section~\ref{sec:coupled-background-convergence}. The representative
of each finite history has exactly its prescribed activities; the
embedding is the same one used in the finite-test estimates. Its
norm is at most $1/8$. At every retained scale its path is
$\widetilde R_{N,m}$, by the construction's input/output identities.
The scalar in its coupling recursion is
$\widetilde\alpha_{N,m}$, not a newly chosen correction.

The finite real-test estimates and their ray consequence in
Section~\ref{sec:background-ray-estimates} give the two one-step
comparison bounds in
\eqref{eq:coupled-activity-step}--\eqref{eq:coupled-cost-step}.
Their constants are $8/128$ for the activity and $4s_m$ for the
scalar correction. The entry-threshold calculation at the end of
Section~\ref{sec:coupled-background-convergence} gives
$\sum_{j<m}s_j\le B_m/128$ for these same allowances. The
larger-width path limit above, the convergent stationary jets and
the retained free horizon supply the background data in the
common-variable comparison. The compensated parameter estimates
of Sections~\ref{sec:local-parameter-comparison} and
\ref{sec:finite-volume-parameter-continuity}, with the suppliers
of Section~\ref{sec:parameter-comparison-suppliers}, then give its
local continuity, including the projected scalar output. Each
application retains the funded transition and output identities
stated there. To apply this continuity to every auxiliary history
with a prescribed coupling script and zero terminal activity, one
must verify those transition hypotheses for its actual input at each
step. In particular, the protected chart, funding and renewed prices,
downward-map bounds and output identification must refer to that input
and the same path and dictionaries. Availability of a one-step map on
the allowed ball alone does not establish these supplier hypotheses.
Their matching for the auxiliary inputs remains an input to the
application here.

When these same-input hypotheses are supplied, all hypotheses of
Section~\ref{sec:coupled-background-convergence} are on one and the
same family. For any terminal sequence
$N\to\infty$ and any corresponding finite trajectories whose bottom
couplings tend to $b$, there are, at every fixed $m,p$, limits
\[
 \beta_{N,m}\longrightarrow\beta_m,\qquad
 x_{N,m}\longrightarrow x_m\ \text{in }\mathcal H_{m,p},\qquad
 q_{N,m}\longrightarrow q_m.
\]
Here $q_{N,m}$ is the scalar activity correction, distinguished from
the path correction $\alpha_{N,m}$. They obey
\[
                \beta_m-\beta_{m+1}-\alpha_m=q_m.
\]
Two sequences with the same bottom limit $b$ give the same limits:
interleave them, whose terminal indices still tend to infinity,
and apply the same uniqueness proof. No compactness of a ball in
$\mathcal H_{m,p}$ is used. The norm convergence and the positive
partition lower bound in Section~\ref{sec:fixed-rung-convergence}
now apply to the actual normalized densities at that scale.

This application retains the strengthened stationary construction,
its convergent free horizon and the funded nonlinear transition.
The argument establishes compatibility of their consequences; the
analytic production of these inputs remains the responsibility of
the corresponding construction estimates.
